\section{Algorithm}\label{sec:algorithm}
Each task is matched to a TC2038 technique that fits its structure, and each main algorithm is compared with an alternative or a baseline.

\subsection{T1: graph traversal}
Connectivity is a reachability question that involves no weights, so graph traversal is sufficient. BFS (with a FIFO queue) and DFS (with an explicit stack) each visit every vertex reachable from the start exactly once; restarting from each vertex not yet visited yields the connected components. Both traversals run on the adjacency list, where the neighbours of a vertex are read directly, and on the adjacency matrix, where the whole row of the vertex must be scanned to find them. The comparison between the two representations is the subject of the task.

\subsection{T2: geometric computation and dynamic programming}
Each pipe length is one Euclidean distance between its endpoints. Network distances between all pairs of nodes come from Floyd-Warshall, a dynamic program over the set of vertices allowed as intermediate points,
\begin{equation}
    D^{(k)}_{ij}=\min\bigl(D^{(k-1)}_{ij},\;D^{(k-1)}_{ik}+D^{(k-1)}_{kj}\bigr),
\end{equation}
which also records a next-hop matrix from which any route can be reconstructed. All $V(V-1)/2$ pairs are needed and $V\le71$, so its $O(V^3)$ cost is small, and the same table later serves as the metric closure of T6. The circuity $d_G(u,v)/\lVert p_u-p_v\rVert\ge1$ measures how much longer the route through the pipes is than the straight line. To query lengths, three structures are compared: a linear scan of the pipe list (baseline), a hash index keyed by the unordered pair of endpoints, and an array of lengths sorted once and searched by bisection, which answers threshold and range queries.

\subsection{T3: greedy shortest paths for sectorization}
All lengths are non-negative, so Dijkstra's greedy rule (permanently settle the unsettled vertex closest to the source, then relax its pipes) yields exact shortest paths. Multi-source Dijkstra places every source in the priority queue at distance 0. This is equivalent to running Dijkstra from a virtual super-source joined to every source by a zero-length pipe: each node is settled by the source that reaches it first, which is exactly $\sigma(v)$, and the whole assignment costs one Dijkstra run. The alternative runs Dijkstra once per source and assigns each node to the source with the smallest distance. Both variants break exact ties in favour of the source listed first. The set $C$ of pipes whose endpoints lie in different sectors is both necessary and sufficient to separate the sectors, since any open inter-sector pipe would join two of them. Closing $C$ does not lengthen any node's route to its own source: if $\sigma(v)=s$, every node on a shortest path from $s$ to $v$ is also closer to $s$ than to any other source, so the whole path stays inside the sector.

\subsection{T4: shortest-path trees for freshness}
T4 reuses the multi-source Dijkstra run of T3 instead of recomputing distances. Dijkstra leaves a parent pointer at every node, forming a shortest-path tree rooted at the sources; the farthest node $t_s$ of each sector is an $O(V)$ scan of the distances within $V_s$, and its route is recovered by following the parent pointers from $t_s$ back to $s$. To compare distances before and after sectorization, one Dijkstra run per source on the open network gives the farthest node that source would have to reach if it supplied the whole network alone (its eccentricity).

\subsection{T5: network flow}
Each pipe of a sector becomes two opposite arcs of capacity $c(u,v)$ in a residual graph, so water may flow in either direction. The Ford-Fulkerson method repeatedly finds a path from $s$ to $t_s$ with spare residual capacity and pushes flow along it until no such path remains. By the max-flow min-cut theorem, the nodes still reachable from $s$ in the final residual graph define a minimum cut: the pipes crossing it are saturated, they are the bottleneck, and their capacities add up to the maximum flow. Choosing augmenting paths arbitrarily gives no polynomial bound for real-valued capacities, so two variants with guarantees are compared. Edmonds-Karp always augments along a shortest path in number of arcs, found by BFS. Dinic builds a BFS level graph and saturates it with a blocking flow found by DFS, using current-arc pointers so that no arc is retried within a phase. Shortest path and maximum flow answer different questions: a shortest path minimises the total length of one route and ignores capacity, while a maximum flow ignores length and asks how much water all routes together can carry. A single route carries at most the capacity of its narrowest pipe, whereas the maximum flow is limited by the narrowest cut, which may span several parallel routes.

\subsection{T6: dynamic programming, greedy, local search and branch and bound}
Let $D$ be the metric closure of the network, $D_{uv}=d_G(u,v)$, computed by Floyd-Warshall. Any feasible sampling walk, read as the order in which it first visits each node, is a tour of $D$ that is no longer than the walk; conversely, any tour of $D$ expands into a walk of the same length by replacing each leg with its shortest path. The optimal sampling walk is therefore the optimal travelling-salesman tour of $D$ starting at the office $o$, expanded back into real pipes with the next-hop matrix. This problem is NP-hard, so no known exact algorithm is polynomial: an exact method is affordable only for small instances, and heuristics are needed for the rest.

Three methods solve the tour problem. Nearest neighbour (greedy) starts at the office and always moves to the closest unvisited node. 2-opt (local search) starts from the nearest-neighbour tour and replaces two legs $(a,b),(c,d)$ by $(a,c),(b,d)$ whenever this shortens the tour, until no such exchange exists; the result, called NN + 2-opt below, is deterministic and is a local optimum, with no guarantee. Branch and bound is exact. It searches partial routes from the office depth first, trying the nearest unvisited nodes first, starts with the NN + 2-opt tour as the incumbent, and discards a partial route ending at $u$ with unvisited set $R$ as soon as
\begin{equation}
    \mathrm{LB}=\mathrm{cost}+\mathrm{MST}(R)+\min_{r\in R}D_{ur}+\min_{r\in R}D_{ro}
\end{equation}
reaches the incumbent's length. Any completion is a path through all of $R$, whose length is at least the minimum spanning tree of $R$, plus one leg from $u$ into $R$ and one leg from $R$ back to $o$, so the bound never overestimates and pruning never discards the optimum. The search is stopped after 60\,s; if it finishes, the incumbent is proven optimal, and otherwise the bound at the root gives a lower limit on the optimum. Only branch and bound guarantees optimality.

\subsection{T7: divide and conquer for geometric search}
Brute force compares the new point $q$ with all $n$ non-source candidates. A 2-d tree divides the candidates recursively at the median of alternating coordinates. A query descends towards $q$ and visits the other side of a split only if the splitting line is closer to $q$ than the best candidate found so far, which prunes most of the tree when the points are well spread. New points are inserted one at a time in file order with the next free id, as node 38 in the slides' FOS example, and each is linked to the nearest non-source node of the network as it stands, so an earlier new node is also a candidate. The third value of each record is taken as the capacity of the new pipe, and its length is the Euclidean distance. Updating the adjacency list adds one key and two entries, while the adjacency matrix cannot grow in place and must be reallocated and copied.

\subsection{Selection hyper-heuristic for T6}\label{sec:hh-design}
A hyper-heuristic searches over heuristics rather than over solutions: a high-level rule chooses, at each step, which low-level heuristic (LLH) to apply, and an acceptance criterion decides whether the result replaces the current solution. A selection hyper-heuristic does this over a fixed set of existing heuristics and can learn online which of them are paying off. Here it works on the same tour problem as T6. The four LLHs are random moves that keep the office first: a 2-opt move (reverse a random segment), a swap of two locations, a relocation of one location, and an or-opt move of a block of 2 or 3 consecutive locations, possibly reversed. Two selection rules are compared. Random selection draws an LLH uniformly at each step; adaptive selection uses probability matching, choosing each LLH with probability proportional to an exponential average (rate 0.05) of whether its recent applications improved the current tour, with a 5\% floor so that no LLH is abandoned. Candidates are accepted by late acceptance hill climbing (LAHC): a candidate is kept if it is no longer than the current tour or than the current tour 50 steps earlier, which lets the search leave local optima without a temperature schedule. Each run applies 30,000 LLHs starting from the nearest-neighbour tour. The baselines are the deterministic NN + 2-opt of T6 and each LLH used alone under the same acceptance and budget, named ``LAHC 2-opt move'', ``LAHC swap move'' and so on.
